\section{Pair spectra as smeared finite-volume densities}
\label{app:pair-spectra}

In finite volume the two-hadron spectrum in each quantum-number sector
is discrete, so the inclusive pair distribution
Eq.~\eqref{eq:pairspec} is defined through a smearing kernel of width
$\sigma$,
\begin{equation}
\rho_\sigma(M) = \sum_n |\langle n|\,\Pi_{h_1h_2}\,|{\rm out}\rangle|^2\,
K_\sigma(M-E_n^{\rm cm}),
\end{equation}
the real-time counterpart of smeared spectral reconstruction in
Euclidean correlators \cite{Hansen:2019idp}.  Three regimes:

\emph{(i) Narrow resonance, fine spectrum} ($\Delta E \ll \Gamma \ll
\sigma$): the line shape is resolved and yields follow from fits with
the same scheme dependence experiment has.

\emph{(ii) $\Gamma \lesssim \Delta E$}: single finite-volume levels
carry the resonance; yields per level map to infinite-volume line
shapes through the L\"uscher formalism
\cite{Luscher:1990ux,Briceno:2017max}, and multi-volume running
reconstructs the distribution.  For the $\Tcc$ benchmark the required
finite-volume technology, including the left-hand-cut subtlety of the
nearby $D^*$ pole, exists
\cite{Padmanath:2022cvl,Raposo:2023oru}.

\emph{(iii) Kernel independence}: ratios of yields extracted with a
common kernel $K_\sigma$ are independent of $\sigma$ up to
$O(\sigma/\Delta M)$ corrections, $\Delta M$ the distance to the nearest
spectral feature --- the statement that makes Tier-3 \emph{ratios}
well-defined even where individual line shapes are scheme-labeled.

The complementarity theorem of Sec.~\ref{sec:resonances} is then
sharp: the fireball out-state supplies the weights
$|\langle n|\Pi|{\rm out}\rangle|^2$ (populations), while Euclidean
spectroscopy supplies the identification of each $|n\rangle$
(identities); the two share one finite-volume spectrum.
